\documentclass[12pt,a4paper]{article}

% ----- Encoding and language -----------------------------------------
\usepackage[utf8]{inputenc}
% \usepackage[T1]{fontenc}
\usepackage{lmodern}
\usepackage[spanish,es-tabla,es-nodecimaldot, es-noquoting]{babel}
\usepackage{microtype}

% ----- Page layout --------------------------------------------------- 
\usepackage[a4paper, left=2.6cm, right=2.6cm, top=2.8cm, bottom=2.6cm,
headheight=15pt]{geometry}
\linespread{1.15}
\setlength{\parskip}{4pt}
\setlength{\parindent}{0.5in}

% ----- Colors --------------------------------------------------------
\usepackage[table]{xcolor}
\definecolor{darkblue}{HTML}{0B2545}
\definecolor{mediumblue}{HTML}{1D4E89}
\definecolor{lightblue}{HTML}{DCE7F3}
\definecolor{borderblue}{HTML}{B8C7DD}
\definecolor{lightgray}{HTML}{F5F6F8}

% ----- Packages ------------------------------------------------------
\usepackage{graphicx,float,booktabs,tabularx,array,multirow,longtable}
\usepackage{enumitem,amsmath,amssymb,caption,xparse,url}
\usepackage{siunitx}
\sisetup{output-decimal-marker={,}, inter-unit-product=\ensuremath{{}\cdot{}}}
\usepackage{tikz}
\usetikzlibrary{arrows.meta,calc,positioning}
\usepackage{pgfplots}
\pgfplotsset{compat=1.16}
\usepackage[most]{tcolorbox}
\usepackage{fancyhdr,titlesec}

% ----- Header and footer ---------------------------------------------
\pagestyle{fancy}
\fancyhf{}
\fancyfoot[L]{\sffamily\footnotesize\bfseries\color{darkblue} Facultad de
Ingeniería Industrial}
\fancyfoot[R]{\sffamily\footnotesize\bfseries\color{darkblue}Pag \thepage}
\renewcommand{\headrulewidth}{0pt}
\renewcommand{\footrulewidth}{0.8pt}
\renewcommand{\footrule}{{\color{darkblue}\hrule width\headwidth
height\footrulewidth}}

% ----- Section titles -----------------------------------------------
\titleformat{\section}
  {\sffamily\bfseries\Large\color{darkblue}}{\thesection.}{0.6em}{}
  [{\vspace{-0.6ex}\color{mediumblue}\titlerule[1.3pt]}]
\titleformat{\subsection}
  {\sffamily\bfseries\large\color{mediumblue}}{\thesubsection}{0.6em}{}
\titleformat{\subsubsection}
  {\sffamily\bfseries\normalsize\color{darkblue}}{\thesubsubsection}{0.6em}{}
\titlespacing*{\section}{0pt}{2.4ex plus 1ex minus .2ex}{1.6ex}
\titlespacing*{\subsection}{0pt}{2.0ex plus .8ex minus .2ex}{1.0ex}
\titlespacing*{\subsubsection}{0pt}{1.6ex plus .5ex}{0.6ex}
\addto\captionsspanish{\renewcommand{\contentsname}{CONTENIDO}}

% ----- Captions and tables ------------------------------------------
\captionsetup{font=small,labelfont={bf,color=darkblue},labelsep=period,
              textfont={color=black!80},skip=6pt}
\renewcommand{\arraystretch}{1.22}
\setlist{itemsep=2pt,topsep=3pt,parsep=0pt}

% Table header cell (white text on dark blue)
\newcommand{\encab}[1]{\textcolor{white}{\sffamily\bfseries #1}}

% ----- Text boxes ----------------------------------------------------
\newtcolorbox{nota}[1][Nota]{colback=lightblue,colframe=mediumblue,coltitle=white,
  fonttitle=\sffamily\bfseries\small,title=#1,arc=2pt,boxrule=0.6pt,
  left=7pt,right=7pt,top=4pt,bottom=4pt,breakable}
\newtcolorbox{enunciado}{colback=lightblue,colframe=darkblue,arc=2pt,boxrule=0pt,
  leftrule=3pt,left=8pt,right=8pt,top=4pt,bottom=4pt,breakable,fontupper=\sffamily\small}

% ----- Image with placeholder box if file does not exist -------------
% Usage \image[width][box height]{file}{description}
\NewDocumentCommand{\image}{O{0.6} O{5cm} m m}{%
  \IfFileExists{#3}{\includegraphics[width=#1\linewidth]{#3}}{%
    {\setlength{\fboxsep}{0pt}\fcolorbox{borderblue}{lightgray}{%
      \parbox[c][#2][c]{#1\linewidth}{\centering\sffamily\small\color{darkblue}%
      \textbf{[ Insert image
]}\\[3pt]\texttt{\detokenize{#3}}\\[3pt]\footnotesize #4}}}}}

% ----- Links ---------------------------------------------------------
\usepackage[colorlinks=true, breaklinks=true, linkcolor=darkblue,
  urlcolor=mediumblue, citecolor=darkblue,
  pdftitle={Laboratorio N\textdegree 3: Ensayo de Corte},
  pdfauthor={Ahon Saenz, Maria Julia; Rey de Castro, Luis Alejandro},
  pdfsubject={Procesos de Manufactura — UNMSM}]{hyperref}
\urlstyle{same}
\makeatletter\g@addto@macro\UrlBreaks{\do\-}\makeatother
\Urlmuskip=0mu plus 2mu\relax
\setlength{\emergencystretch}{2.5em}
\newcommand{\refbib}[1]{[\ref{#1}]}
\newcommand{\pendiente}[1]{\colorbox{yellow!60}{\textbf{[#1]}}}

% =====================================================================
\begin{document}
% =====================================================================

% ===== COVER PAGE ====================================================
\begin{titlepage}
	\thispagestyle{empty}
	\begin{tikzpicture}[remember picture,overlay]
		\fill[darkblue] (current page.north west) rectangle ($(current page.north
			east)+(0,-4.4cm)$);
		\fill[mediumblue] ($(current page.north west)+(0,-4.4cm)$) rectangle
		($(current page.north east)+(0,-4.65cm)$);
		\fill[darkblue] (current page.south west) rectangle ($(current page.south
			east)+(0,1.1cm)$);
		\node[text=white,align=center,text width=16cm,font=\sffamily] at ($(current
			page.north)+(0,-2.25cm)$) {%
			{\Large\bfseries UNIVERSIDAD NACIONAL MAYOR DE SAN MARCOS}\\[5pt]
			{\small\itshape Universidad del Perú, Decana de América}\\[10pt]
			{\normalsize\bfseries ESCUELA PROFESIONAL DE INGENIERÍA INDUSTRIAL}};
	\end{tikzpicture}

	\vspace*{2cm}

	\begin{center}
		\image[0.22][4.3cm]{img/unmsm-logo.png}{}

		\vspace{0.8cm}
		{\sffamily\bfseries\color{darkblue}\fontsize{22}{26}\selectfont
			Laboratorio N\textdegree 3: Ensayo de Corte\par}

		\vspace{0.4cm}

		\begin{tcolorbox}[colback=white,colframe=white,arc=3pt,width=0.9\linewidth,
				box align=center,left=12pt,right=12pt,top=8pt,bottom=8pt]
			\centering\sffamily
			\begin{tabular}{@{}>{\bfseries\color{darkblue}}r@{\hspace{1em}}l@{}}
				Asignatura: & Procesos de Manufactura (L1-3) \\[5pt]
				Docente:    & Mg. Ing. Víctor Rosales Urbano \\
			\end{tabular}
		\end{tcolorbox}

		\vspace{0.4cm}
		{\sffamily\bfseries\color{darkblue}\large Integrantes\par}
		\vspace{5pt}
		\sffamily\begin{tabular}{@{}l@{\hspace{2.5em}}r@{}}
			Ahon Saenz, Maria Julia Armida      & (23170001) \\
			Rey de Castro Rupay, Luis Alejandro & (22170163) \\
		\end{tabular}

		\vspace{1cm}

		\image[0.3][5.2cm]{img/maju.jpeg}{}
		\hspace{1.5cm}
		\image[0.3][5.2cm]{img/luis.jpg}{}
	\end{center}

	\vfill
	\begin{center}
		{\sffamily\bfseries\color{darkblue}\large Lima, 2026\par}
	\end{center}
\end{titlepage}
\setcounter{page}{2}

% ===== CONTENIDO =====================================================
\tableofcontents
\newpage

% ===== 1. PRESENTACIÓN ===============================================
\section{PRESENTACIÓN}

El mecanizado por arranque de viruta es uno de los grupos de procesos más
importantes de la manufactura: partiendo de una pieza en bruto, una herramienta
de corte remueve el material sobrante hasta obtener la forma, las dimensiones y
el acabado que pide el plano. Aunque suele asociarse con máquinas como el torno
o la fresadora, el principio físico es el mismo cuando la herramienta se mueve a
mano: un filo penetra en el material, lo deforma por corte y lo separa en forma
de viruta. Entender ese mecanismo es lo que permite, después, elegir bien la
herramienta, la velocidad y el fluido de corte.

\vspace{6pt}

En este tercer laboratorio realizamos una operación de aserrado manual:
cortamos una pieza cilíndrica de hierro (Fe) con un arco de sierra,
sujetándola en un tornillo de banco. El docente marcó previamente la pieza para
que la longitud resultante fuera de aproximadamente \SI{70}{\milli\metre}. Los
datos de la práctica (pieza, herramienta, tiempos e inspección) se anotaron en
la hoja de registro del laboratorio, cuya información se presenta en la sección
de resultados.

\vspace{6pt}

La pieza cortada no es un producto final, sino el material de partida de
la siguiente operación: en la hoja de registro, la operación de destino es el
torno. Esto ilustra una idea central de la manufactura: cada operación
recibe una pieza de la anterior y le deja un sobrante (sobreespesor) para
que la siguiente pueda trabajar. Por eso, aun siendo una operación sencilla, el
aserrado se registra con el mismo rigor que una operación de máquina.

\vspace{6pt}

\begin{figure}[H]
	\centering
	\image[0.75][4.6cm]{img/presentacion.jpeg}{Foto general de la experiencia de
		laboratorio}
	\caption{Sistema físico general de la experiencia de laboratorio}
\end{figure}

% ===== 2. MARCO TEÓRICO ==============================================
\section{MARCO TEÓRICO}

\subsection{Teoría de corte y mecanizado}

La teoría de corte estudia los fenómenos que ocurren cuando una
herramienta entra en contacto con un material y produce su separación mediante
deformación y desprendimiento. Durante este proceso intervienen factores como la
geometría de la herramienta, las propiedades del material, las fuerzas de corte,
la fricción y las condiciones de trabajo (Trent \& Wright, 2000).

\vspace{6pt}

El mecanizado es un proceso de manufactura en el que una herramienta de
corte se utiliza para remover el exceso de material de una pieza, de modo que
lo que queda tenga la forma deseada. La acción principal consiste en aplicar
una deformación por corte para formar la viruta y exponer una nueva
superficie (Groover, 2007).

\vspace{6pt}

Los procesos de remoción de material se agrupan en tres categorías
(Tabla~\ref{tab:categorias}). El aserrado pertenece a la primera: es un proceso
de mecanizado convencional en el que una herramienta de múltiples filos, la hoja
de sierra, remueve material en forma de viruta.

\vspace{6pt}

\begin{table}[!htb]
	\centering
	\caption{Categorías básicas de los procesos de remoción de
		material}\label{tab:categorias}
	\small
	\rowcolors{2}{white}{lightblue}
	\begin{tabularx}{\linewidth}{l X}
		\rowcolor{darkblue}
		\encab{Categoría}       & \encab{Descripción y ejemplos}
		\\
		Mecanizado convencional & Una herramienta de corte afilada remueve material
		en forma de viruta:
		torneado, fresado, taladrado, limado y aserrado.
		\\
		Procesos abrasivos      & Partículas abrasivas duras remueven material:
		rectificado, lapeado,
		pulido.
		\\
		Procesos no tradicionales
		                        & Usan energía química, eléctrica, térmica o
		mecánica sin herramienta
		de corte convencional: electroerosión (EDM), láser, chorro de agua.
		\\
	\end{tabularx}
\end{table}

\vspace{2pt}

En la práctica realizada, estos principios se aplicaron mediante el aserrado
manual de una barra maciza de acero de sección circular. La hoja de sierra fue
desplazada de manera repetitiva sobre una línea previamente trazada hasta
conseguir la separación del trozo requerido.

\subsection{Mecánica del corte de metales}

En el corte de metales, la herramienta penetra en el material y lo deforma
plásticamente a lo largo de un plano de corte (o de cizallamiento). El
material que sale por ese plano forma la viruta, que se desliza sobre la cara de
ataque de la herramienta. El modelo más usado para analizarlo es el
corte ortogonal, en el que el filo es perpendicular a la velocidad de
corte y el problema puede tratarse en dos dimensiones. Cuando el filo está
inclinado respecto de esa dirección se habla de corte oblicuo, que es
tridimensional. La Figura~\ref{fig:ortogonal} resume la geometría del corte
ortogonal.

\vspace{6pt}

\begin{figure}[H]
	\centering
	\begin{tikzpicture}[scale=0.9,>=Stealth]
		\coordinate (A) at (5,0);
		\coordinate (B) at (2.86,1.5);
		% pieza de trabajo
		\filldraw[fill=black!15,draw=darkblue,line width=0.8pt]
		(0,-1.2) -- (9.5,-1.2) -- (9.5,0) -- (A) -- (B) -- (0,1.5) -- cycle;
		% viruta
		\filldraw[fill=lightblue,draw=darkblue,line width=0.8pt]
		(A) -- (B) -- (3.53,4.01) -- (5.67,2.51) -- cycle;
		% herramienta
		\filldraw[fill=mediumblue!35,draw=darkblue,line width=0.8pt]
		(A) -- (5.88,3.28) -- (8.6,3.28) -- (8.6,0.5) -- cycle;
		% líneas auxiliares
		\draw[dashed,mediumblue] (0,0) -- (A);
		\draw[dashed,mediumblue] (A) -- (5,3.7);
		% ángulos
		\draw[darkblue] ($(A)+(180:1.3)$) arc (180:145:1.3);
		\node[font=\small,darkblue] at ($(A)+(162:1.75)$) {$\phi$};
		\draw[darkblue] ($(A)+(90:1.7)$) arc (90:75:1.7);
		\node[font=\scriptsize,darkblue] at ($(A)+(82.5:2.05)$) {$\alpha$};
		% espesores
		\draw[<->,darkblue] (1.2,0) -- (1.2,1.5)
		node[midway,left,font=\small] {$t_0$};
		\draw[<->,darkblue] (5.336,1.256) -- (2.963,1.892)
		node[midway,above,sloped,font=\scriptsize] {$t_c$};
		% velocidad
		\draw[->,very thick,darkblue] (9.4,3.8) -- (7.6,3.8)
		node[midway,above,font=\small] {$V$};
		% rótulos
		\node[font=\sffamily\scriptsize,darkblue] at (7.3,2.2) {Herramienta};
		\node[font=\sffamily\scriptsize,darkblue] at (4.35,3.0) {Viruta};
		\node[font=\sffamily\scriptsize,darkblue] at (7.5,-0.7) {Pieza de trabajo};
		\node[font=\sffamily\scriptsize,darkblue,align=center] (pc) at (1.3,2.9)
		{Plano de\\corte};
		\draw[->,darkblue,thin] (1.6,2.35) -- (3.6,1.0);
	\end{tikzpicture}
	\caption{Esquema del corte ortogonal: $t_0$ es el espesor antes del corte,
		$t_c$ el espesor de la viruta, $\phi$ el ángulo del plano de corte,
		$\alpha$ el ángulo de ataque y $V$ la velocidad de la
		herramienta.}\label{fig:ortogonal}
\end{figure}

\subsection{Aserrado manual}

El aserrado es un proceso de corte en el que una herramienta con múltiples
dientes elimina material de manera progresiva. Cada diente funciona como un
pequeño filo que penetra en la pieza y desprende una cantidad de material,
formando una ranura de corte. La acción sucesiva de los dientes permite
profundizar esta ranura hasta completar la separación de la pieza (Choudhury,
2021).

\vspace{6pt}

Para acero, las hojas suelen ser de acero rápido o bimetálicas. El parámetro
principal de selección es el número de dientes por pulgada (TPI): como criterio
práctico, debe haber al menos dos o tres dientes en contacto con el material
en todo momento; por ello, para barra maciza se usan pasos más gruesos (del
orden de 18 TPI) y para tubos o perfiles delgados, más finos (24 a 32 TPI). La
pieza se sujeta en un tornillo de banco lo más cerca posible de la línea de
corte, para reducir la vibración y evitar que la hoja se atasque o se rompa.

\vspace{6pt}

En el aserrado manual, el movimiento de corte es proporcionado por el operario
mediante carreras de avance y retroceso. Por ello, la técnica utilizada, la
presión aplicada y la velocidad de las carreras pueden influir en el rendimiento
de la operación.

\subsection{Formación de viruta}

Una viruta es un fragmento de material residual con forma de lámina curvada o
espiral que es extraído mediante un cepillo u otras herramientas. Su forma
depende del material, de la geometría de la herramienta y de las condiciones de
corte, y se clasifica como se muestra en la Tabla~\ref{tab:virutas}.

\vspace{6pt}

La formación de viruta ocurre cuando el material situado delante del filo de
corte se deforma hasta separarse de la pieza. En términos generales, el material
experimenta una deformación plástica localizada y posteriormente se desprende
debido a la acción de la herramienta.

\vspace{6pt}

\begin{table}[!htb]
	\centering
	\caption{Clasificación de las virutas}\label{tab:virutas}
	\small
	\rowcolors{2}{white}{lightblue}
	\begin{tabularx}{\linewidth}{>{\raggedright\arraybackslash}p{3.2cm} X}
		\rowcolor{darkblue}
		\encab{Según su\ldots} & \encab{Tipos}
		\\
		Tipo                   & Discontinua; continua; continua con protuberancias.
		\\
		Clase                  & Plástica; cortada; de arranque.
		\\
		Forma                  & Agujas; desmenuzadas; bastoncitos; trozos espirales o
		helicoidales; espirales netas; trozos cortos de cinta; hélices cortas o
		largas,
		estrechas o anchas; virutas de sesgo rectilíneo.
		\\
	\end{tabularx}
\end{table}

\vspace{2pt}

En el caso de la sierra manual, cada diente remueve una pequeña cantidad de
material durante la carrera de corte. La repetición de este proceso genera
progresivamente la ranura que permite separar la pieza. Por ello, aunque el
corte se realice manualmente, se presentan los mismos fenómenos básicos de
remoción de material estudiados en la teoría de corte (Choudhury, 2021).

\subsection{Herramientas de corte y desgaste}

Las herramientas de corte se clasifican según varios criterios
(Tabla~\ref{tab:clasif}). La herramienta utilizada en este laboratorio fue una
hoja bimetálica para metales montada en un arco de sierra. Sus dientes
constituyen los elementos encargados de penetrar en el acero y producir el
desprendimiento del material.

\vspace{6pt}

\begin{table}[!htb]
	\centering
	\caption{Clasificación de las herramientas de corte}\label{tab:clasif}
	\small
	\rowcolors{2}{white}{lightblue}
	\begin{tabularx}{\linewidth}{>{\raggedright\arraybackslash}p{4.3cm} X}
		\rowcolor{darkblue}
		\encab{Criterio}           & \encab{Clases}
		\\
		Según el número de filos   & De un filo; de doble filo o en hélice; de filos
		múltiples.
		\\
		Según el material          & Acero de herramientas (WS); aceros de
		herramientas aleados (SS);
		metales duros aleados (HS); diamantes.
		\\
		Por el movimiento de corte & Fijo; contra el material; en contra dirección.
		\\
		Por el tipo de viruta      & Viruta continua; en forma de coma; polvo sin
		forma definida.
		\\
	\end{tabularx}
\end{table}

\vspace{2pt}

Durante el corte, los dientes están sometidos a esfuerzos y fricción. Como
consecuencia, puede producirse desgaste progresivo de la herramienta, lo que
reduce su capacidad de corte y puede aumentar el esfuerzo requerido para avanzar
sobre el material (Trent \& Wright, 2000). Por esta razón, antes de iniciar una
operación de aserrado es importante
comprobar que la hoja se encuentre correctamente tensada y que sus dientes estén
en condiciones adecuadas.

\subsection{Condiciones de corte}

Las condiciones de corte son los parámetros que el operario o el programador
fija para una operación. Los conceptos principales que intervienen en el proceso
son:

\vspace{2pt}

\begin{itemize}
	\item \textbf{Metal sobrante (sobreespesor):} diferencia entre el espesor
	      inicial $E$ y el final $e$ de la pieza, $T=E-e$ (mm).
	\item \textbf{Profundidad de corte:} espesor de material removido en una
	      pasada. En el cilindrado, $t=\dfrac{D_f-D_i}{2}$, donde $D_i$ y $D_f$
	      son
	      los diámetros inicial y final (mm).
	\item \textbf{Velocidad de avance:} desplazamiento de la herramienta respecto
	      de la pieza por unidad de giro, de tiempo o de carrera.
	\item \textbf{Velocidad de corte:} velocidad relativa entre el filo y el
	      material en el punto más desfavorable, $V_c=\pi D n$ (m/min o ft/min),
	      con $D$ el diámetro correspondiente (en metros) y $n$ las revoluciones
	      por
	      minuto de la pieza o de la herramienta.
\end{itemize}

Para producir el corte es necesario aplicar una fuerza que permita que los
dientes de la hoja penetren en el material. Esta fuerza depende de factores como
las propiedades del acero, la geometría de los dientes, la cantidad de material
que se remueve y las condiciones de contacto entre la herramienta y la pieza.

\vspace{6pt}

En el aserrado manual, la fuerza es proporcionada directamente por el operario.
Por ello, aplicar la presión principalmente durante la carrera activa y mantener
un ritmo constante permite aprovechar mejor el movimiento de la herramienta. Los
estudios de mecánica del corte consideran las fuerzas generadas durante la
interacción entre herramienta y pieza como uno de los aspectos fundamentales
para analizar el proceso de remoción de material (Choudhury, 2021).

\subsection{Sujeción y precisión del corte}

La sujeción de la pieza es fundamental para mantenerla estable durante el
proceso. Una pieza correctamente fijada disminuye los movimientos y facilita que
la hoja siga la línea previamente trazada.

\vspace{6pt}

\begin{table}[!htb]
	\centering
	\caption{Principales elementos que intervienen en el aserrado manual}
	\small
	\rowcolors{2}{white}{lightblue}
	\begin{tabularx}{\linewidth}{
			>{\raggedright\arraybackslash}l
			>{\raggedright\arraybackslash}X
		}
		\rowcolor{darkblue}
		\encab{Elemento}  &
		\encab{Función en el proceso}
		\\
		Pieza de trabajo  &
		Material sobre el cual se realiza el corte.
		\\
		Hoja de sierra    &
		Penetra en el material y remueve pequeñas cantidades de acero mediante sus
		dientes.
		\\
		Arco de sierra    &
		Permite sujetar y tensar la hoja y transmitir el movimiento.
		\\
		Tornillo de banco &
		Mantiene inmovilizada la pieza durante el corte.
		\\
		Línea de trazado  &
		Sirve como referencia para controlar la dirección del corte.
		\\
	\end{tabularx}
\end{table}

\vspace{2pt}

La precisión del corte es importante debido a que la pieza será utilizada
posteriormente en una operación de torneado. Por ello, después del aserrado se
realizó el desbarbado y la verificación de la pieza antes de considerarla
habilitada.

% ===== 3. MATERIALES Y EQUIPOS =======================================
\section{MATERIALES Y EQUIPOS}

La práctica consistió en habilitar un trozo de barra maciza de acero de sección
circular mediante aserrado manual, utilizando un arco de sierra y un tornillo de
banco para mantener fija la pieza. El objetivo fue obtener un trozo de
aproximadamente \SI{70}{\milli\metre} de longitud, considerando el material
adicional necesario para el posterior mecanizado en torno. La operación se
realizó completamente de forma manual sobre el banco de trabajo, sin utilizar
una máquina de corte. Para controlar las dimensiones y la dirección del corte se
emplearon instrumentos de medición y trazado.


\subsection{Pieza de trabajo}

El material utilizado fue una barra maciza de acero (Fe) de sección circular,
procedente del almacén del taller, identificada en la hoja de registro como
pieza N\textdegree 01. La barra tenía una sección de \SI{1}{in} y
una longitud inicial de \SI{6000}{\milli\metre}. De ella se habilitó un trozo de
aproximadamente 70 mm de longitud.

\begin{figure}[H]
	\centering
	\image[0.5][4.6cm]{img/pieza_marcada.jpeg}{FOTO 1: pieza cilíndrica de Fe}
	\caption{Pieza de Fe con la marca de corte trazada por el docente.}
\end{figure}

\subsection{Herramientas de corte y sujeción}

Para realizar el corte se utilizó un arco de sierra para metales, equipado con
una hoja bimetálica de 300 mm (12"). El arco permitió mantener la hoja tensada y
transmitir el movimiento generado por el operario. La barra fue fijada mediante
un tornillo de banco o morsa mecánica, evitando que se desplazara durante el
corte.

\begin{table}[!htb]
	\centering
	\caption{Herramientas utilizadas durante la práctica}
	\small
	\rowcolors{2}{white}{lightblue}
	\begin{tabularx}{\linewidth}{
			>{\raggedright\arraybackslash}l
			>{\raggedright\arraybackslash}X
			>{\raggedright\arraybackslash}X
		}
		\rowcolor{darkblue}
		\encab{Herramienta}                                         &
		\encab{Especificación}                                      &
		\encab{Función}
		\\
		Arco de sierra para metales                                 &
		Bastidor regulable, mango tipo pistola y tensor de mariposa &
		Sujetar y tensar la hoja durante el corte.
		\\
		Hoja de sierra bimetálica                                   &
		300 mm (12")                                                &
		Realizar el corte y remover material mediante sus dientes.
		\\
		Tornillo de banco                                           &
		Mordazas paralelas, fijado al banco                         &
		Sujetar e inmovilizar la barra.
		\\
	\end{tabularx}
\end{table}

\begin{figure}[H]
	\centering
	\image[0.65][4.6cm]{img/arco_sierra.jpeg}{FOTO 2: arco de sierra con su hoja}
	\caption{Arco de sierra y hoja utilizada en el laboratorio.}
\end{figure}

\begin{figure}[H]
	\centering
	\image[0.85][4.6cm]{img/tornillo_banco.jpeg}{FOTO 3: tornillo de banco con la
		pieza sujeta}
	\caption{Tornillo de banco con la pieza sujeta para el aserrado.}
\end{figure}

\subsection{Instrumentos de trazado y verificación}
Antes de realizar el corte fue necesario determinar la longitud de la pieza y
marcar la zona donde se realizaría el aserrado. Para ello se utilizó un vernier
o pie de rey, un rayador de punta de carburo y una escuadra de precisión
(realizado por el docente).

\vspace{6pt}

El vernier permitió medir la pieza y establecer la longitud requerida. El
rayador permitió marcar la línea de corte sobre las caras de la barra, mientras
que la escuadra facilitó el trazado perpendicular y posteriormente la
verificación de la superficie obtenida.

\begin{table}[!htb]
	\centering
	\caption{Instrumentos de trazado y verificación}
	\small
	\rowcolors{2}{white}{lightblue}
	\begin{tabularx}{\linewidth}{
			>{\raggedright\arraybackslash}l
			>{\raggedright\arraybackslash}X
			>{\raggedright\arraybackslash}X
		}
		\rowcolor{darkblue}
		\encab{Instrumento}                                               &
		\encab{Especificación}                                            &
		\encab{Función}
		\\
		Vernier o pie de rey                                              &
		Rango 0–150 mm, resolución 0,05 mm                                &
		Medir la sección y determinar la longitud de corte
		\\
		Tiza color anaranjado                                             &
		Roca sedimentaria compuesta principalmente de carbonato de calcio &
		Marcar la línea de corte sobre la barra
		\\
	\end{tabularx}
\end{table}

\begin{figure}[H]
	\centering
	\image[0.45][4.6cm]{img/instrumentos.png}{Instrumentos utilizados para el
		trazado y verificación}
	\caption{Instrumentos utilizados para el trazado y
		verificación.}\label{fig:instrumentos}
\end{figure}

\subsection{Equipos de protección personal}

Debido a que la operación fue realizada manualmente y existía contacto con una
herramienta de corte, aristas metálicas y rebabas, se utilizaron equipos de
protección personal para reducir los riesgos durante la práctica.

\begin{table}[H]
	\centering
	\caption{Equipos de protección personal utilizados}
	\small
	\rowcolors{2}{white}{lightblue}
	\begin{tabularx}{\linewidth}{
			>{\raggedright\arraybackslash}l
			>{\raggedright\arraybackslash}X
		}
		\rowcolor{darkblue}
		\encab{EPP}                &
		\encab{Riesgo que controla}
		\\
		Guantes de seguridad       &
		Cortes, aristas vivas y contacto con rebabas.
		\\
		Botas con puntera de acero &
		Caída de la barra o de la pieza cortada.
		\\
		Mandil o guardapolvo       &
		Contacto con virutas, grasa y suciedad.
		\\
	\end{tabularx}
\end{table}

% ===== 4. PROCEDIMIENTO ==============================================
\section{PROCEDIMIENTO}

La secuencia de trabajo se resume en la Figura~\ref{fig:flujo} y se detalla a
continuación.

\begin{enumerate}[label=\textbf{\arabic*.}]
	\item Se completaron los datos generales de la hoja de registro: pieza,
	      práctica, material, operación, herramienta, tiempo estimado y
	      operaciones de origen y destino.
	\item El docente marcó sobre la pieza de Fe la longitud de corte, de
	      aproximadamente \SI{70}{\milli\metre}.
	\item Se fijó la pieza en el tornillo de banco, con la marca de corte
	      accesible para la hoja y lo más cerca posible de las mordazas
	      (Figura~\ref{fig:proceso}).
	\item Se anotó la hora de inicio (17:58) y se realizó el corte con el arco
	      de sierra siguiendo la marca, hasta separar el tramo de la pieza.
	\item Al terminar se anotó la hora final (18:22), se calculó el tiempo
	      total de la operación y se realizó la medición de longitud y peso.
	\item El docente inspeccionó la pieza y firmó el campo de inspección de la
	      hoja de registro.
\end{enumerate}

\vspace{6pt}

\begin{figure}[H]
	\centering
	\begin{tikzpicture}[node distance=5mm,
			caja/.style={rectangle,rounded
					corners=3pt,fill=darkblue,text=white,font=\sffamily\scriptsize\bfseries,
					minimum height=1.15cm,minimum width=2.3cm,align=center,inner sep=3pt},
			flecha/.style={-{Stealth[length=2.2mm]},thick,color=mediumblue}]
		\node[caja] (a) {Recibir pieza\\marcada};
		\node[caja,right=of a] (b) {Sujetar en\\tornillo de banco};
		\node[caja,right=of b] (c) {Aserrar con\\arco de sierra};
		\node[caja,right=of c] (d) {Registrar\\resultados};
		\node[caja,right=of d] (e) {Inspección\\del docente};
		\draw[flecha] (a)--(b); \draw[flecha] (b)--(c); \draw[flecha] (c)--(d);
		\draw[flecha] (d)--(e);
	\end{tikzpicture}
	\caption{Secuencia del procedimiento experimental.}\label{fig:flujo}
\end{figure}

\vspace{6pt}

\begin{figure}[H]
	\centering
	\image[0.65][4.6cm]{img/proceso_aserrado.jpeg}{FOTO 4: operación de aserrado
		en curso (pieza en el tornillo de banco y arco de sierra)}
	\caption{Aserrado manual de la pieza de Fe.}\label{fig:proceso}
\end{figure}

% ===== 5. ELABORACIÓN DE RESULTADOS ==================================
\section{ELABORACIÓN DE RESULTADOS}\label{sec:resultados}

\subsection{Hoja de registro de la práctica}

La Figura~\ref{fig:hoja} muestra la hoja de registro llenada durante el
laboratorio, y la Tabla~\ref{tab:hoja} transcribe sus datos.

\begin{figure}[H]
	\centering
	\image[0.85][6cm]{img/hoja_datos.png}{Hoja de registro de la práctica}
	\caption{Hoja de registro de la práctica N\textdegree 03 (original
		manuscrita).}\label{fig:hoja}
\end{figure}

\begin{table}[!htb]
	\centering
	\caption{Datos de la hoja de registro (transcripción)}\label{tab:hoja}
	\small
	\rowcolors{2}{white}{lightblue}
	\begin{tabularx}{\linewidth}{>{\bfseries\raggedright\arraybackslash}p{5.2cm}
			X}
		\rowcolor{darkblue}
		\encab{Campo}            & \encab{Dato
			                           registrado}
		\\
		Horario                  & Lab 03 (Lunes 17:20 - 19:00)
		\\
		Fecha                    & 21/09/2026
		\\
		N\textdegree{} de pieza  & 01
		\\
		Descripción del artículo & Fe, Ø \SI{1}{in} $\times$ \SI{69}{\milli\metre}
		\\
		Práctica / fecha         & 03 \ -- \ 21/09/2026
		\\
		Tipo de material         & Fe
		\\
		Peso final               & Approx. \SI{340}{g}
		\\
		Operación actual         & Aserrado
		\\
		Máquina                  & -- (operación manual, sin máquina)
		\\
		Herramienta              & Arco con hoja de sierra
		\\
		Tiempo estimado          & \SI{30}{\minute}
		\\
		Operación origen         & Aserrado
		\\
		Operación destino        & Torno
		\\
		Tiempo de inicio         & 17:58
		\\
		Tiempo final             & 18:22
		\\
		Tiempo total             & \SI{24}{\minute}
		\\
		Inspección               & Mg. Ing. Víctor Rosales Urbano
		\\
	\end{tabularx}
\end{table}

\begin{figure}[H]
	\centering
	\image[0.55][4.6cm]{img/pieza_cortada.jpeg}{FOTO 5: tramo de Fe cortado (y,
		si se conserva, la viruta o limalla generada)}
	\caption{Pieza cortada al finalizar la operación.}
\end{figure}

\subsection{Análisis de tiempos}

El tiempo total anotado (\SI{24}{\minute}) es consistente con las horas de
inicio y fin: $18{:}22 - 17{:}58 = \SI{24}{\minute}$. Comparado con el tiempo
estimado de \SI{30}{\minute}, la operación se completó la operación terminó
dentro de lo previsto (Tabla~\ref{tab:tiempos}, Figura~\ref{fig:tiempos}).
El tiempo real fue un 20\,\% menor que el estimado. Sin embargo, con una sola
pieza no es posible generalizar este resultado: depende de la destreza del
operario, del estado de la hoja y del diámetro de la pieza.

\begin{table}[!htb]
	\centering
	\caption{Tiempo estimado frente a tiempo real}\label{tab:tiempos}
	\small
	\rowcolors{2}{white}{lightblue}
	\begin{tabular}{l r}
		\rowcolor{darkblue}
		\encab{Indicador}                 & \encab{Valor}      \\
		Tiempo estimado                   & \SI{30}{\minute}   \\
		Tiempo real (18:22 $-$ 17:58)     & \SI{24}{\minute}   \\
		Diferencia (real $-$ estimado)    & $-\SI{6}{\minute}$ \\
		Tiempo real / tiempo estimado     & 80\,\%             \\
		Reducción respecto de lo estimado & 20\,\%             \\
	\end{tabular}
\end{table}

\begin{figure}[H]
	\centering
	\begin{tikzpicture}
		\begin{axis}[
				ybar, width=0.62\linewidth, height=5.6cm,
				symbolic x coords={Estimado,Real}, xtick=data,
				enlarge x limits=0.55, ymin=0, ymax=36,
				ytick={0,10,20,30}, ylabel={Tiempo (min)},
				bar width=30pt, ymajorgrids, grid style={borderblue!70},
				nodes near coords, nodes near coords style={font=\sffamily\small},
				label style={font=\sffamily\small},
				tick label style={font=\sffamily\small},
				axis x line*=bottom, axis y line*=left]
			\addplot[fill=mediumblue,draw=darkblue] coordinates {(Estimado,30)
					(Real,24)};
		\end{axis}
	\end{tikzpicture}
	\caption{Tiempo estimado y tiempo real de la operación de
		aserrado.}\label{fig:tiempos}
\end{figure}

% ===== 6. CUESTIONARIO ===============================================
\section{CUESTIONARIO}

\subsection*{1. ¿Qué elementos de mecanizado existen?}

En toda operación de mecanizado intervienen los siguientes elementos:

\begin{itemize}
	\item \textbf{Máquina-herramienta:} proporciona el movimiento y la energía del
	      corte (torno, fresadora, taladro, cepillo, sierra mecánica). En nuestro
	      laboratorio no hubo máquina: la energía la puso el operario.\\[-10pt]
	\item \textbf{Herramienta de corte:} el elemento con filo que remueve el
	      material (cuchilla, fresa, broca, hoja de sierra). Se define por su
	      material, su geometría y su número de filos.\\[-10pt]
	\item \textbf{Pieza de trabajo:} el material a mecanizar, con su forma inicial
	      y el metal sobrante que debe removerse.\\[-10pt]
	\item \textbf{Dispositivos de sujeción:} mantienen la pieza (o la herramienta)
	      en posición: mandril, plato, prensa o tornillo de banco.\\[-10pt]
	\item \textbf{Movimientos de la operación:} movimiento principal o de corte,
	      movimiento de avance y movimiento de profundidad o penetración.\\[-10pt]
	\item \textbf{Fluido de corte:} lubricante o refrigerante que reduce la
	      fricción y el calor en la zona de corte.\\[-10pt]
	\item \textbf{Viruta:} el material removido, resultado de la operación.
\end{itemize}

\subsection*{2. ¿Qué tipos de cortes existen?}

Según los criterios de la sesión, los cortes pueden clasificarse de la siguiente
manera:

\begin{itemize}
	\item \textbf{Por la geometría del corte:} \emph{ortogonal}, en el que el filo
	      es perpendicular a la velocidad de corte y el problema es
	      bidimensional; y \emph{oblicuo}, en el que el filo está inclinado y el
	      problema es tridimensional.\\[-10pt]
	\item \textbf{Por el número de filos de la herramienta:} de un filo
	      (cuchilla de torno), de doble filo o en hélice (broca) y de filos
	      múltiples (fresa, hoja de sierra).\\[-10pt]
	\item \textbf{Por el tipo de viruta:} continua, continua con protuberancias
	      (filo recrecido) y discontinua.\\[-10pt]
	\item \textbf{Por la clase de viruta:} plástica, cortada y de
	      arranque.\\[-10pt]
	\item \textbf{Por la operación:} torneado, fresado, taladrado, aserrado, entre
	      otras. El aserrado realizado en este laboratorio es un corte con una
	      herramienta de filos múltiples y movimiento alternativo.
\end{itemize}

\subsection*{3. ¿Cuáles son los parámetros de corte?}

Los parámetros de corte son la velocidad de corte, el avance y la profundidad de
corte (Tabla~\ref{tab:parametros}). Junto con el metal sobrante, son los
conceptos principales del proceso según la sesión.

\begin{table}[!htb]
	\centering
	\caption{Parámetros de corte}\label{tab:parametros}
	\small
	\rowcolors{2}{white}{lightblue}
	\begin{tabularx}{\linewidth}{>{\raggedright\arraybackslash}p{3.6cm}
			>{\raggedright\arraybackslash}p{3.6cm} X}
		\rowcolor{darkblue}
		\encab{Parámetro}    & \encab{Expresión / unidad}        &
		\encab{Significado}
		\\
		Velocidad de corte   & $V_c=\pi D n$ (m/min)             & Velocidad
		relativa entre el filo y
		el material en el punto más desfavorable ($D$ en m; $n$ en rpm).
		\\
		Avance               & mm/rev, mm/diente o mm/carrera    & Desplazamiento de
		la herramienta
		respecto de la pieza por unidad de giro o de carrera.
		\\
		Profundidad de corte & mm; en cilindrado $t=(D_f-D_i)/2$ & Espesor de
		material removido
		en una pasada.
		\\
		Metal sobrante       & $T=E-e$ (mm)                      & Material total
		que debe
		removerse entre la pieza en bruto y la pieza final.
		\\
	\end{tabularx}
\end{table}

\subsection*{4. ¿Cuáles son los componentes principales para el corte?}

El corte requiere, como mínimo, una \emph{herramienta}, una \emph{pieza} y un
\emph{movimiento relativo} entre ambas. Los componentes de la herramienta y del
sistema que se estudian son los siguientes:

\begin{itemize}
	\item \textbf{Superficie de desprendimiento o de ataque:} cara sobre la que
	      se desliza la viruta.\\[-10pt]
	\item \textbf{Filo principal y filo secundario de corte:} aristas que separan
	      el material.\\[-10pt]
	\item \textbf{Caras de incidencia principal y secundaria:} caras que
	      enfrentan la superficie recién cortada; su ángulo evita el
	      roce.\\[-10pt]
	\item \textbf{Punta de la herramienta:} unión de los dos filos.\\[-10pt]
	\item \textbf{Mango y plaquita:} cuerpo de la herramienta y el inserto de
	      corte
	      que se fija a él.\\[-10pt]
	\item \textbf{Movimiento primario (de corte), movimiento de avance y
		      profundidad,} que definen la trayectoria de la herramienta.\\[-10pt]
	\item \textbf{Viruta:} material removido que sale por la cara de ataque.
\end{itemize}

\subsection*{5. ¿Cuáles son las tres categorías básicas de procesos de remoción
	de material?}

Son el \textbf{mecanizado convencional} (corte con herramienta de filo, como
torneado, fresado, taladrado y aserrado), los \textbf{procesos abrasivos}
(rectificado, lapeado, pulido) y los \textbf{procesos no tradicionales}
(electroerosión, láser, chorro de agua y otros que remueven material sin una
herramienta de corte convencional). Ver Tabla~\ref{tab:categorias}.

\subsection*{6. Indique algunas razones por las que el mecanizado es comercial y
	tecnológicamente importante}

Según Groover, el mecanizado es importante porque:

\begin{itemize}
	\item \textbf{Es aplicable a una gran variedad de materiales:} metales,
	      plásticos y cerámicos, con las herramientas adecuadas.\\[-10pt]
	\item \textbf{Puede generar casi cualquier forma geométrica:} superficies
	      planas, cilíndricas, agujeros, roscas y formas complejas.\\[-10pt]
	\item \textbf{Alcanza tolerancias estrechas:} típicamente del orden de
	      $\pm\SI{0.025}{\milli\metre}$ o mejores, superiores a las de la mayoría
	      de
	      procesos de formado.\\[-10pt]
	\item \textbf{Produce buenos acabados superficiales:} rugosidades del orden de
	      $R_a\approx\SI{0.4}{\micro\metre}$ o menores.\\[-10pt]
	\item \textbf{Complementa a otros procesos:} piezas fundidas o forjadas se
	      mecanizan después para lograr agujeros roscados, superficies de ajuste
	      y otras características de precisión.\\[-10pt]
	\item \textbf{Es económico en lotes pequeños:} no requiere moldes ni matrices,
	      por lo que es adecuado para prototipos y series cortas.
\end{itemize}

Como contrapartida, genera desperdicio de material en forma de viruta y suele
requerir más tiempo que los procesos de formado.

\subsection*{7. ¿Cuáles son los parámetros de una operación de mecanizado que se
	incluyen dentro del alcance de las condiciones de corte?}

Las condiciones de corte de una operación comprenden la velocidad de
corte $V_c$, el avance $f$, la profundidad de corte y el uso (o no) de un fluido
de corte y su tipo.

\vspace{6pt}

Estos valores se eligen según el material de la pieza, el material de la
herramienta y el objetivo de la pasada (desbaste o acabado). En nuestro aserrado
manual, la velocidad quedó determinada por el ritmo del operario, el avance por
la fuerza aplicada y la profundidad por el progreso de la ranura; ninguno de
estos valores se midió.

% ===== 7. CONCLUSIONES ===============================================
\section{CONCLUSIONES}

\begin{itemize}
	\item Se realizó el corte de la pieza N\textdegree 01 de Fe con arco de sierra
	      y tornillo de banco, respetando la marca del docente de aproximadamente
	      \SI{70}{\milli\metre}, y se registró la operación en la hoja del
	      laboratorio con herramienta, tiempos e inspección.\\[-10pt]
	\item La operación tomó \SI{24}{\minute} frente a los \SI{30}{\minute}
	      estimados, es decir, el 80\,\% del tiempo previsto. Con una sola pieza
	      no puede generalizarse este resultado, pero indica que la estimación
	      fue razonable.\\[-10pt]
	\item En el aserrado manual las condiciones de corte (velocidad, avance y
	      profundidad) las determina el operario y no pueden fijarse con precisión
	      como en una máquina. Esto hace que el tiempo y la calidad del corte
	      dependan de la destreza y del estado de la hoja.\\[-10pt]
	\item Cada diente de la hoja actúa como una herramienta de un filo que remueve
	      una cantidad mínima de material; por eso se necesitan muchas carreras
	      para atravesar la pieza. El mecanismo es el mismo que en el corte
	      ortogonal estudiado en la teoría.\\[-10pt]
	\item La pieza cortada queda lista como material de partida para el torno, lo
	      que muestra la lógica de encadenar operaciones dejando un sobrante para
	      la siguiente.
\end{itemize}

% ===== 8. RECOMENDACIONES ============================================
\section{RECOMENDACIONES}

\begin{itemize}
	\item Registrar el tipo de hoja (dientes por pulgada) y, si es posible, contar
	      las carreras por minuto durante un intervalo corto, para poder estimar la
	      velocidad de corte.\\[-10pt]
	\item Elegir el paso de la hoja según la sección de la pieza: al menos dos o
	      tres dientes deben estar en contacto con el material, y más finos si se
	      trata de un tubo o perfil delgado.\\[-10pt]
	\item Sujetar la pieza firmemente y cerca de la línea de corte; cortar con
	      presión en la carrera de avance y aliviarla en el retorno; iniciar con
	      una
	      pequeña ranura guía sobre la marca.\\[-10pt]
	\item Usar lentes de seguridad y retirar las rebabas del corte con una lima
	      antes de manipular o medir la pieza.
\end{itemize}

% ===== 9. ENSAYOS VIRTUALES ==========================================
\section{ENSAYOS VIRTUALES}

No se encontró un simulador gratuito específico del aserrado manual. Los
siguientes recursos permiten reforzar la teoría del corte y los ensayos con
fuerzas de corte que no pudieron realizarse en esta práctica:

\begin{enumerate}[label=\textbf{\arabic*.}]
	\item \textbf{Virtual Labs (Ministerio de Educación de la India).} Portal de
	      laboratorios virtuales de ingeniería, con experimentos web de
	      manufactura y procesos de formado y mecanizado. \\
	      \url{https://www.vlab.co.in}\\[-4pt]
	\item \textbf{Monitoring of Cutting Forces in Turning} (Sredanovic y Carou,
	      2024). Capítulo de acceso abierto que describe un procedimiento de
	      laboratorio para medir las fuerzas de corte en el torneado y analizar su
	      efecto en la precisión dimensional. \\
	      \url{https://doi.org/10.1007/978-3-031-48468-1_8}\\[-4pt]
	\item \textbf{Cutting forces during turning (GUNT).} Descripción de un equipo
	      didáctico con transductor de tres componentes para medir fuerza de
	      corte,
	      de avance y pasiva, equivalente al dinamómetro citado en la teoría. \\
	      \url{https://www.gunt.de/en/products/mechatronics/manufacturing-engineering/technological-experiments/cutting-forces-during-turning/054.10200/ft102/glct-1:pa-148:ca-67:pr-367}
\end{enumerate}

% ===== 10. ACCIDENTES RELACIONADOS ===================================
\section{ACCIDENTES RELACIONADOS}

La pieza cortada pasa al torno, y los accidentes más graves en talleres de
mecanizado ocurren con piezas en rotación. Se mencionan tres casos documentados
por organismos oficiales:

\begin{itemize}
	\item \textbf{Maquinista atrapado por la manga en un torno (Oregón,
		      EE.\,UU.).}
	      Un maquinista de 69 años con más de 30 de experiencia murió cuando la
	      manga de su ropa se enredó en la pieza que giraba en un torno. Un
	      compañero pulsó el botón rojo de la máquina para detenerla, pero no era
	      un paro de emergencia y la máquina siguió funcionando. El informe
	      recomienda una política de vestimenta para operadores, formación en
	      peligros y supervisión de conductas seguras. \\
	      \url{https://www.cdc.gov/niosh/face/pdfs/12or018.pdf}\\[-4pt]
	\item \textbf{Torno con protecciones anuladas (Wisconsin, 2016).} Un operador
	      de 36 años se enredó en el husillo mientras pulía a mano un cilindro
	      metálico de 40 pulgadas y murió dos días después. La OSHA halló que la
	      empresa dejaba operar el torno con los enclavamientos de seguridad
	      anulados y la sancionó por una infracción deliberada y una grave. \\
	      \url{https://www.dol.gov/newsroom/releases/osha/osha20160906}\\[-4pt]
	\item \textbf{Guante atrapado en un chavetero (Bethlehem Steel, inspección de
		      1991).} Un
	      maquinista con 12 años de experiencia pulía con lija el cuello de un
	      rodillo que giraba a 344 rpm y usaba guantes de cuero; el puño de un
	      guante quedó atrapado en un chavetero y murió. Existía una norma escrita
	      que prohibía los guantes y había sido advertido esa mañana. \\
	      \url{https://www.osha.gov/ords/imis/accidentsearch.accident_detail?id=889444}
\end{itemize}

En el aserrado manual los riesgos son distintos pero también reales: cortes por
la hoja o por rebabas, rotura de la hoja y desplazamiento de una pieza mal
sujeta.

% ===== 11. BIBLIOGRAFÍA ==============================================
\section{BIBLIOGRAFÍA}

\begin{enumerate}[label={[\arabic*]},ref=\arabic*,leftmargin=2.2em,itemsep=2pt]
	\item Choudhury, A. R. (2021). \emph{Metal cutting and machine tools.} NPTEL,
	      Indian Institute of Technology Kharagpur. NPTEL – Metal Cutting and
	      Machine Tools\\[-8pt]
	\item Trent, E. M., \& Wright, P. K. (2000). \emph{Metal cutting} (4th ed.).
	      Butterworth-Heinemann. Metal Cutting – ScienceDirect\\[-8pt]
	\item Groover, M. P. (2007). \emph{Fundamentos de
		      manufactura moderna: materiales, procesos y sistemas} (3\textdegree
	      ed.). McGraw-Hill.\\[-8pt]
	\item Kalpakjian, S., \& Schmid, S. R. (2014).
	      \emph{Manufactura, ingeniería y tecnología} (7\textdegree ed.). Pearson
	      Educación.\\[-8pt]
	\item Oregon FACE Program. (2014). \emph{Experienced
		      Journeyman Machinist Killed While Operating an Engine Lathe} (Informe
	      OR 2012-18-1). Oregon Institute of Occupational Health Sciences.
	      \url{https://www.cdc.gov/niosh/face/pdfs/12or018.pdf}\\[-8pt]
	\item U.S. Department of Labor, OSHA. (2016). \emph{OSHA's
		      investigation of lathe operator's fatal injuries finds machine guards
		      were bypassed at Wisconsin machining facility} [Comunicado de prensa].
	      \url{https://www.dol.gov/newsroom/releases/osha/osha20160906}\\[-8pt]
	\item U.S. Department of Labor, OSHA. (s.\,f.).
	      \emph{Accident
		      Report Detail, Summary Nr. 889444: Employee killed when head strikes
		      against lathe.}
	      \url{https://www.osha.gov/ords/imis/accidentsearch.accident_detail?id=889444}\\[-8pt]
	\item Virtual Labs, Ministry of Education (India).
	      \url{https://www.vlab.co.in}\\[-8pt]
	\item Sredanovic, B., \& Carou, D. (2024). Monitoring of
	      Cutting Forces in Turning. En \emph{Materials Forming, Machining and
		      Tribology} (pp. 131--149). Springer.
	      \url{https://doi.org/10.1007/978-3-031-48468-1_8}
\end{enumerate}

\end{document}
